%%==========================================================
%%  CAPÍTULO 7: CONCLUSIONES
%%==========================================================

Las conclusiones deben ser la respuesta a los objetivos o propositos planteados así como presentar las lecciones aprendidas más importantes durante el proceso. Este capitulo debería incluir: 

\begin{itemize}
\item \textbf{Resumen de los hallazgos}: Las conclusiones deben proporcionar un resumen claro y conciso de los principales hallazgos de la investigación. Esto debería incluir una revisión de los objetivos de la investigación y una discusión de cómo se cumplieron estos objetivos.
\item \textbf{Relación con la literatura existente}: Las conclusiones deben relacionar los hallazgos de la investigación con la literatura existente. Esto podría implicar la comparación de los resultados con estudios previos, la discusión de cómo los hallazgos contribuyen a la teoría existente, etc.
\item \textbf{Implicaciones prácticas}: Si es relevante, las conclusiones deben discutir las implicaciones prácticas de los hallazgos de la investigación. Esto podría incluir recomendaciones para la práctica profesional, sugerencias para la implementación de políticas, etc.
\item \textbf{Limitaciones de la investigación}: Las conclusiones deben reconocer las limitaciones de la investigación. Esto podría incluir limitaciones en la metodología, limitaciones en la recopilación de datos, etc.
\item \textbf{Sugerencias para futuras investigaciones}: Las conclusiones deben proporcionar sugerencias para futuras investigaciones. Esto podría incluir áreas que no se exploraron en la investigación actual, preguntas que surgieron como resultado de la investigación, etc.
\item \textbf{Claridad y coherencia}: Las conclusiones deben estar bien escritas y presentadas de manera clara y coherente. Deben proporcionar un cierre efectivo para la tesis.
\end{itemize}

\section{Conclusiones}

\section{Trabajos futuros}

\section{Lecciones aprendidas}

\section{Puede agregar más secciones si lo necesita}